\section{Model and the Data-augmentation Approach}
\label{sec:model}

In this section, we establish the theoretical framework. We focus on conjoint analysis as the primary use case to present the theory and conduct the empirical analysis detailed in Section \ref{sec:empirical-1}.

\subsection{Setup}
\label{sec:setup}

We first provide a detailed explanation of the setup, starting with an in-depth discussion of the theoretical framework underlying the data generation process. We will outline the assumptions, structure, and specific parameters that guide this process, building a foundation for understanding how the data is conceptualized and modeled within this framework.

\subsubsection{Data Generation Process and the Best-in-class Estimation.}
Consider a setting where we observe two datasets. The first, referred to as the \textit{primary data}, consists of $m$ data points, 
$\left\{\big(\x^{\sf P}_{i}, y^{\sf P}_{i}, z^{\sf P}_{i}\big)\right\}_{i = 1}^m$.
For a fixed positive integer $k$, we denote the set $\{1,\dots, k\}$ by $\Kscr$ and define $\Kscr^+ = \Kscr \cup \{0\}$. 
\revise{These data points are independently and identically distributed (i.i.d.) according to the same distribution as the random vector $(\x, y, z)$, where $y, z \in \Kscr^+$ and $\x \in \Real^{dk}$.
To clarify, in the context of conjoint analysis with AI-generated labels, $\x$ represents the features or context of a choice setting, and $y$ denotes the respondent's choice. Here, $\Kscr^+$ corresponds to the $k+1$ options available to the respondent, where $0$ represents the outside option. Additionally, $\x = (\x^{\top}_{(1)}, \x^{\top}_{(2)}, \dots, \x^{\top}_{(k)})^{\top}$, where $\x^{\top}_{(j)} \in \Real^d$ corresponds to the features of option $j$ for each $j \in \Kscr$.} Throughout this paper, we refer to the random variable $y$ as the \textit{real label}. 
The variable $z$ represents a predicted label generated by an AI model based on $\x$, and we refer to it as the \textit{AI-generated label}. 
The second dataset, called the \textit{auxiliary data}, consists of $\left\{\big(\x_{i}, z_{i}\big)\right\}_{i = 1}^n$, which are independent of the primary data but follow the same distribution. However, in the auxiliary data, the true label $y_i$ is missing for each data point. 
In Figure \ref{fig:illustration-data}, we illustrate this theoretical setup using an example simplified from the empirical setting in Section \ref{sec:empirical-1}. 
\revise{In contrast, in typical conjoint analysis settings, the AI-generated labels $z$ is not available, and 
our goal is to fit the vector $\betavec  \in \Real^{d}$ to capture the choice probabilities as accurate as possible using  Multinomial Logit (MNL) models, i.e., 
$\betavec$ such that 
\begin{align}
\label{eqn:mnl-conjoint}
    \sigma_j(
\x;\betavec 
)~=~ \frac{e^{\x^{\top}_{(j)}\betavec}}{1+\sum_{\ell \in \Kscr} e^{\x^{\top}_{(\ell)}\betavec}},~~\forall\, j \in \Kscr, 
\end{align}
best approximates $\Prob(y\,|\,x)$.
}


How is this data generation process implemented in practice? A manager facing a conjoint study problem can begin by collecting a small sample of human response data. Using the same questions presented to human respondents, the manager can query an LLM to generate the primary dataset containing both labels $y$ and $z$. Subsequently, the manager can generate additional responses by querying the LLM with a different set of choices, which forms the auxiliary dataset where the true label $y$ is missing. Since the cost of querying an LLM is negligible compared to recruiting human subjects in many cases, the auxiliary data is typically much larger than the primary data in size.



\begin{figure}[h!]
    \centering
    \includegraphics[width=165mm]{fig/illustration-data-cropped.pdf}
    \caption{Illustration of the dataset}
    \label{fig:illustration-data}
    \begin{minipage}{\textwidth}
    \footnotesize
    \begin{tabular}{p{1cm} p{15cm}}
         \hspace{-1mm}\textbf{Notes}: & \setlength{\baselineskip}{1.2\baselineskip} In this figure, we illustrate the two datasets, using an example simplified from the empirical setting in Section \ref{sec:empirical-1}. 
         For each data point in the primary data of size $m$, two vaccine options are presented to both the human respondent and the AI. 
         Corresponding to the first data point, the first vaccine has the following features: 50\% efficacy, a one-year protection duration, a major side-effect rate of one in a million, a minor side-effect rate of one in ten, and endorsement by President Donald Trump. These features can be converted into a numerical feature vector, $\x^{\sf P}_{1, (1)}$, suitable for conjoint analysis, for example, using one-hot encoding. 
         Note that the options chosen by the human subject and AI may agree with each other (e.g., in the last data point) or differ from each other (e.g., in the first data point). 
         The auxiliary data, which is of size $n$, is different from the primary data in that we only have the label solicited from the AI. 
    \end{tabular}
    
    % \begin{itemize}[leftmargin=0pt]
    %     \item [\textbf{Notes}:] In this figure, we illustrate the primary dataset using an example simplified from the empirical setting in Section \ref{sec:empirical-1}. For instance, corresponding to the first data point, two vaccine options are presented to the human respondent. The first vaccine has the following features: 50\% efficacy, a one-year protection duration, a major side-effect rate of one in a million, a minor side-effect rate of one in ten, and endorsement by President Donald Trump. These features can be converted into a numerical feature vector, $\x^{\sf P}_{1, (1)}$, suitable for conjoint analysis, for example, using one-hot encoding. Based on these features, the AI recommends selecting the first vaccine (i.e., $z^{\sf P}_{1} = 1$), which aligns with the human respondent’s choice (i.e., $y^{\sf P}_{1} = 1$). Note that the second human respondent opts to select no vaccine from the two offered. The auxiliary dataset is similar to the primary dataset, except that it lacks the human-generated labels for these data points.
    % \end{itemize}
  \end{minipage}
\end{figure}
We assume that the distribution $\Prob(y\,|\,\x, z)$ can be parameterized as $\Prob(y = j\,|\,\x, z)\,=\,g_j(\x, z; \thetavec^*)$ for each $j \in \Kscr^+$, where $\thetavec^* \in \Real^q$ is not known to us. 
This assumption is motivated by the fact that modern artificial AI models, such as an LLM, are trained with vast amounts of text data, including information from numerous sources on the internet, that may contain product reviews, messaging boards, and other online forums with contributions from a wide range of consumers discussing the products they shop for and purchase \citep{brand2023using}.
Therefore, AI can ``mimic" how respondents make choices, 
leading to highly informative AI-generated labels. Therefore, the model $\big(g_j(\x, z; \thetavec^*):j\in \Kscr^+\big)$, which can be chosen as the MNL model or a neural network, 
is specified to capture the conditional probabilities $\Prob(y\,|\,\x, z)$ in practical scenarios.
We also assume throughout that $|g_j(\cdot)| \leq 1$ for all 
$j\in \Kscr^+$, which is reasonable for a wide range of models. 

\revise{
In this work, we frame the main goal of our conjoint analysis as identifying the \textit{best-in-class} estimator of \eqref{eqn:best-in-class}—the one that most closely approximates the true choice probabilities within the set of all MNL models.}\footnote{Note that we do not assume that the data is generated by a MNL model. We recognize that choice probabilities parameterized by the MNL model may be misspecified. In practice, customer choice probabilities can exhibit considerable complexity, yet researchers and practitioners often rely on the simpler MNL model as an approximation in conjoint analyses.} The quality of this approximation is measured by the Kullback–Leibler (KL) divergence. Formally, the best-in-class estimator is defined as follows:
% For each fixed $\betavec \in \Real^{d}$, we use the Kullback–Leibler divergence (KL-divergence) to measure the quality of the MNL model parameterized by $\betavec$. 
% We define 
\begin{align}
\label{eqn:best-in-class}
    \betavec^* ~\in~
    \argmin_{\betavec \in \Real^{d}}
    \left\{
    \E_{\x}\Big[
    {\sf KL}\big( \Prob(y| \x)\,\big\vert\,\sigma_{y}(\x, \betavec) \big)
    \Big]~=~
    \E_{\x}\left[
    \sum_{j\in \Kscr^+} \Prob\big(y = j| \x\big)
    \log\left(
    \frac{\Prob\big(y = j| \x\big)}{\sigma_{j}(\x; \betavec)}
    \right)
    \right]\right\}.
\end{align}
That is, measuring the loss of approximation using the KL-divergence,  
$\big(\sigma_j(\x; \betavec^*):j\in \Kscr^+\big)$ gives us the approximation to the true choice probabilities that is best in the class of MNL models.  
Under reasonable assumptions, this minimizer exists and is uniquely defined (see Theorem \ref{thm:consistency} below). 

We further clarify the distinction between our problem and the measurement error issues discussed in the econometrics literature \citep[e.g., see][]{bound2001measurement} and in recent work by \cite{connell2024estimating}. Specifically, in our context, conditional on $\x$, $z$ should not be interpreted as a noisy realization of $y$, as is commonly assumed in the measurement error framework. In our setting, $y \perp z \,|\, \x$ because the AI cannot access the internal thought and decision-making processes of a human subject, and any relationship between $y$ and $z$ is mediated entirely through $\x$. Nonetheless, we emphasize that our estimation procedure and theoretical results impose no assumptions on the joint distribution of $y$ and $z$ given $\x$. Therefore, even if $z$ indeed represents $y$ measured with error, our approach remains valid.


\subsection{Primary, Auxiliary and Na\"ive Estimators} 
\label{sec:3-estimators}
With no access to auxiliary data such as AI-generated labels, the default approach is to fit the parameter $\betavec$ using only the primary data with its real labels $\left\{\big(\x^{\sf P}_{i}, y^{\sf P}_{i}\big)\right\}_{i = 1}^m$ using the standard Maximum Likelihood Estimation (MLE). We use $\hat{\betavec}^{\sf P}$ to denote the estimator obtained using only the primary data. When the primary data size is sufficiently large, $\hat{\betavec}^{\sf P}$ can be close to $\betavec^*$. However, the size of the primary data is often restricted due to the costs of recruiting real subjects. A small primary set yields a \textit{inaccurate} estimator that is far away from $\betavec^*$. 
On the flipping side, we may directly perform MLE on the auxiliary data $\left\{\big(\x_{i}, z_{i}\big)\right\}_{i = 1}^n$, using $z_{i}$ as labels. 
Let us define this estimator as $\hat{\betavec}^{\sf A}$. Given the distribution of LLM-generated data and human data are different (i.e., the distributions of $ y $ and $ z $, conditioned on $ \x $, are different), it is easy to recognize that $\hat{\betavec}^{\sf A}$ could be severely biased regardless of the size of auxiliary dataset compared to $\betavec^*$.

Due to this limitation, people intend to utilize auxiliary data and AI-generated labels to facilitate the estimation of $\betavec^*$. A natural idea for utilizing the AI labels is to simply pool the primary and auxiliary sets. In particular, we define
$$
(\tilde{\x}_{i}, \tilde{y}_i)~=~
\begin{cases}
    ({\x}_{i}^{\sf P}, y_{i}^{\sf P})~&\text{if } 1 \leq i \leq m \\
    ({\x}_{i-m}, z_{i-m})~&\text{if } m+1 \leq i \leq m+n
\end{cases},
$$
and perform the standard MLE on $(\tilde{\x}_{i}, \tilde{y}_i)_{i=1}^{m+n}$.
We call this approach the \textit{na\"ive augmentation} and write the estimator as $\hat{\betavec}^{\sf Naive}$. The auxiliary data largely enlarges the data size. However, since $z$ and $y$ in general come from different distributions, $\hat{\betavec}^{\sf Naive}$ again can be biased, even with large data size. Note that both $\hat{\betavec}^{\sf P}$ and $\hat{\betavec}^{\sf A}$ are both special cases of $\hat{\betavec}^{\sf Naive}$ by setting $n = 0$ or $m = 0$. In the following, we present a simple example to demonstrate the limitations of these estimators discussed above and summarize the observation in Proposition~\ref{prop:bias}.

\begin{examp}[AI-Generated Data Cannot Replace Human Data]
\label{examp:ai-bias}
{\rm 
Consider a simple setting with 
only one product and no features.
Conditional on $z = j$, where $j \in \Kscr^+ = \{0,1\}$, $y = j$ with probability $p$.
Suppose that $\Prob(z = 1) = \alpha$. 
For clarity, we display the joint distribution in Table \ref{tab:joint-y-z}.

\begin{table}
\renewcommand{\arraystretch}{1.5}
\centering
{\footnotesize
\caption{Joint distribution of $y$ and $z$ in Example \ref{examp:ai-bias}}
\label{tab:joint-y-z}
\begin{tabular}{  c |c  c }
\hline
\hline
  & $y=1$ & $y=0$ \\ \hline
 $z=1~$ & $~\alpha p~$ & $~\alpha (1-p)~$ \\  \hline
 $z=0~$ & $~(1-\alpha) p~$ & $~(1-\alpha) (1-p)~$
 \\\hline\hline
\end{tabular}
}
\end{table}


Therefore, the probability that $y = 1$ is given by $c_1(\alpha, p):=\alpha p +(1 - \alpha)(1 - p)$ and clearly the best-in-class estimator in \eqref{eqn:mnl-conjoint} is given by $\beta^* = \log\left(\frac{c_1(\alpha, p)}{1- c_1(\alpha, p)}\right)$, which is computed through
$
\frac{e^{\beta^*}}{1 + e^{\beta^*}} = \Prob(y = 1) = c_1(\alpha, p)
$.
Suppose that $n, m \rightarrow \infty$ and $n/m \rightarrow \rho \neq 0$. 
Asymptotically, the na\"ive estimator satisfies 
$$
c_2(\alpha, p):=\frac{\alpha \rho}{1 + \rho}+ 
\frac{\alpha p +(1 - \alpha)(1 - p)}{1 + \rho}~\approx~\frac{e^{\hat{\beta}^{\sf Naive}}}{1 + e^{\hat{\beta}^{\sf Naive}}},
$$
or equivalently 
$
\hat{\beta}^{\sf Naive}~\rightarrow~ \log\left(
\frac{c_2(\alpha, p)}{1-c_2(\alpha, p)}
\right).
$
One can verify that unless $\alpha = 1/2$ or $p = 1$,  this limit is not equal to $\beta^*$. 
This is always the case even when $n,m\rightarrow \infty$, as long as $\rho \neq 0$. 
Also note that even when $p$ is close to, but different from, one, $z$ is very informative of $y$ but the na\"ive methods still cannot lead to desired estimates. }
\end{examp}


\begin{prop}[Bias of $\hat{\betavec}^{\sf A}$ and $\hat{\betavec}^{\sf Naive}$.]
\label{prop:bias}
When $m, n \rightarrow \infty$, in general, 
    $\hat{\betavec}^{\sf A}$ and $\hat{\betavec}^{\sf Naive}$ are not consistent estimators of ${\betavec}^*$.
\end{prop}






\subsection{Estimation with AI-augmented Data}
\label{sec:estimation}
These challenges lead to our research question: \textit{How do we extract value from the AI-augmented data to facilitate the estimation of $\betavec^*$?} To answer this question, we propose the following data augmentation approach that allows us to use the AI-generated data to fit the model.\footnote{For simplicity, here we assume that an optimizer $\hat{\betavec}^{\sf AAE}$ is obtained. 
Otherwise, we may optimize approximately and alternatively choose 
$\hat{\betavec}^{\sf AAE}$ such that $
\widehat{Q}(\hat{\thetavec};\hat{\betavec}^{\sf AAE}) = \sup_{\betavec \in \Real^{d}} \widehat{Q}(\hat{\thetavec};\betavec) - o_{\text{P}}(1)
$. 
In this case, our key theoretical result, Theorem \ref{thm:consistency}, is still valid with a slight adaption of the analysis. Interested readers may refer to the proof of Theorem \ref{thm:consistency} in the appendix.}

\vspace{10pt}

\begin{mdframed} \footnotesize \underline{{\sc Estimation with AI-Augmented Data}} 

\noindent \textbf{Step 1.} 
Obtain an estimator $\hat{\thetavec}$ to $\thetavec^*$, where $\Prob(y = j\,|\,\x, z)\,=\,g_j(\x, z; \thetavec^*)$, using the primary data.

\noindent \textbf{Step 2.} With the auxiliary data, we construct the estimator $\hat{\betavec}^{\sf AAE}$ as 
\begin{align}
\label{eqn:AAE-def}
    \hat{\betavec}^{\sf AAE} ~\in~\argmax_{\betavec \in \Real^{d}}\,\left\{ \widehat{Q}(\hat{\thetavec};\betavec)\,=\, \frac{1}{n}\sum_{i=1}^n\sum_{j\in \Kscr^+}g_j(\x_i, z_i; \hat{\thetavec})\log \sigma_j(\x; \betavec)\right\}.
\end{align}
\end{mdframed}

\vspace{10pt}

We refer to $\hat{\betavec}^{\sf AAE}$ as the \textit{AI-augmented estimator} (AAE). 
In the next section, we demonstrate that our estimator $ \hat{\betavec}^{\sf AAE} $ successfully recovers $ \betavec^* $ and exhibits the desired asymptotic properties. The key idea in our approach is to use the primary data---which contains both LLM-generated and human-produced data---to learn an efficient mapping between human and LLM-generated data. We then leverage this mapping to construct the estimator. This strategy closely resembles a transfer learning approach, where knowledge from one domain (in our case, the LLM-generated data) is transferred to another domain (the human-produced data). The fundamental assumption underpinning such transfer learning techniques is that learning the mapping function between two domains, denoted by $g(\cdot)$ in our context, is inherently simpler and requires fewer data points than carrying out the individual learning tasks within either domain on its own. In other words, in our context, the AI model undertakes the ``heavy lifting'' of capturing human choice behavior in a fuzzing way, and the first stage in our estimation is trying to capture how fuzzy this AI model is compared to the human model. Moreover, the mapping will become easier to learn with lower LLM costs and higher LLM quality. This is because, as the accuracy and the size of AI-generated labels improves, we anticipate a high-quality estimation of $ \thetavec$ due to a decrease in the asymptotic variance of $ \thetavec$ estimation; see Section \ref{sec:AAE} and Appendix \ref{app:comparison} for further discussions.

\revise{
We highlight the intuition of our method from two alternative angles. 
First, our approach is motivated by principles from transfer learning and knowledge distillation, though adapted to the context of conjoint analysis. In transfer learning, knowledge from one domain is leveraged to reduce data needs in another \citep{pan2009survey,zhuang2020comprehensive}; similarly, our method transfers information contained in the joint distribution of product features and LLM-generated labels ($\gvec(\x, z; \thetavec^*)$) to improve inference on human preferences. Likewise, while classical knowledge distillation involves a student model mimicking a teacher \citep{hinton2015distilling}, our framework echoes this idea by learning the conditional probability of real labels given LLM labels and product features. Although our method does not involve parameter reuse across tasks or a teacher–student architecture in the narrow sense, the underlying principle is the same: using a large, complex source of information (the LLM) to guide a simpler, statistically grounded estimator.
Such a principle can be used to integrate AI-generated labels in other settings, e.g., a continuous target in a regression problem (see discussions in Section \ref{sec:additional}
and Appendix \ref{app:regression}).
}

\revise{Second, while LLMs encode vast information, their outputs are systematically different from human responses. Importantly, this discrepancy is often easier to learn from small samples than estimating preferences directly. Our method builds on this intuition: we use a small human dataset to capture the gap between human and LLM labels, and then apply this correction to the much larger LLM-generated dataset. In the AAE algorithm, the first step models the discrepancy between AI and human labels conditional on product features $\x$. The second step then adjusts estimation using the predicted human label based on the AI label and features. When the first-stage model is sufficiently accurate, the second-stage estimator reduces variance by leveraging the enlarged AI-labeled dataset (Proposition \ref{prop:weak-dom}). Thus, the more precise the first-stage predictions, the greater the bias reduction in the second stage, as formalized in Theorem \ref{thm:aae-bound}.}


